\documentclass[11pt,a4paper]{article}
\usepackage{fontspec,amsmath,amssymb,amsthm,microtype}
\setmainfont{lmroman10-regular.otf}[BoldFont=lmroman10-bold.otf,ItalicFont=lmroman10-italic.otf,BoldItalicFont=lmroman10-bolditalic.otf]
\setmonofont{lmmono10-regular.otf}[Scale=MatchLowercase]
\usepackage[margin=25mm]{geometry}
\usepackage[unicode,hidelinks]{hyperref}
\setlength{\emergencystretch}{3em}
\setlength{\parskip}{0.4em}
\setlength{\parindent}{0pt}
\allowdisplaybreaks[2]
\raggedbottom
\title{\texorpdfstring{Exact Gibbs invariance of an explicit corrected Verlet proposal}{Exact Gibbs invariance of an explicit corrected Verlet proposal}}
\author{LangevinLean mathematical exposition}
\date{}
\begin{document}\maketitle
\noindent Source module: \texttt{Kinetic\allowbreak{}Langevin/\allowbreak{}Hamiltonian\allowbreak{}Metropolis.\allowbreak{}lean}. The accompanying HTML pairs this mathematical exposition with the exact Lean statements and records the scope of the formalization.


A volume-preserving involution has a symmetric old/\allowbreak{}new-state graph measure. Weighting that graph by a positive target density and applying the ordinary Metropolis ratio makes its accepted flow symmetric. The rejected mass stays at the old state. This module proves exact invariance for the explicit reversed-Verlet proposal. Passing from these corrected proposals to the unadjusted Hamiltonian flow and then to the full kinetic diffusion still requires quantitative limits.


\section{\texorpdfstring{Construct the proposal graph and derive its swap symmetry}{Construct the proposal graph and derive its swap symmetry}}


For \(J^2=\mathrm{id}\), the graph map is \(G(z)=(z,Jz)\). Swapping the old and new state gives \(\operatorname{swap}\circ G=G\circ J\). If \(J\) preserves a base measure \(\lambda\), pushforward composition yields \(\operatorname{swap}_\#G_\#\lambda=G_\#J_\#\lambda=G_\#\lambda\). Lean proves the pointwise graph identity using the involution equation, verifies graph measurability, and rewrites the two pushforward compositions. Thus symmetry of the graph measure is derived from the actual proposal rather than introduced as a detailed-balance field.



{\raggedright\noindent Lean declarations: \texttt{proposal\allowbreak{}Graph}, \texttt{proposal\allowbreak{}Graph\_\allowbreak{}swap\_\allowbreak{}invariant}.\par}

\section{\texorpdfstring{Apply the accepted-\allowbreak{}flow argument to an arbitrary positive density}{Apply the accepted-flow argument to an arbitrary positive density}}


Let the target be \(w\lambda\), with measurable \(w>0\). Its proposal joint law is \(G_\#(w\lambda)\). The pushforward-with-density theorem expresses this as the symmetric base graph weighted by \(w\) at its first coordinate. The Metropolis ratio is \(\min\{1,w(y)/w(x)\}\), so accepted joint density is \(\min\{w(x),w(y)\}\), symmetric under swapping. The checked shared Metropolis theorem proves that accepted new-state mass plus rejected old-state mass has the initial marginal. Lean computes that marginal by projecting the graph back to its first coordinate, giving the exact identity with \(w\lambda\). No invariant measure or acceptance-balance conclusion is assumed.



{\raggedright\noindent Lean declarations: \texttt{metropolis\_\allowbreak{}involution\_\allowbreak{}preserves\_\allowbreak{}density}.\par}

\section{\texorpdfstring{Bound rejection by an actual energy increase}{Bound rejection by an actual energy increase}}


For energy increment \(d\), rejection is \(1-\min\{1,e^{-d}\}\). If \(d\le0\), the proposed energy decreases and rejection is zero. If \(d\ge0\), the exponential tangent inequality \(e^{-d}\ge1-d\) gives rejection at most \(d\). Thus rejection is at most \(\max\{d,0\}\le|d|\). Lean splits the sign cases, rewrites the minimum and maximum with the corresponding exponential bounds, and applies the tangent inequality. This is a scalar probability bound; converting it to an L\ensuremath{{}^{2}} trajectory error also requires the size of the rejected jump and must not silently assume that jump is of order \(h\).



{\raggedright\noindent Lean declarations: \texttt{metropolis\_\allowbreak{}rejection\_\allowbreak{}le\_\allowbreak{}energy\_\allowbreak{}increase}, \texttt{metropolis\_\allowbreak{}rejection\_\allowbreak{}le\_\allowbreak{}abs\_\allowbreak{}energy}.\par}

\section{\texorpdfstring{Instantiate the exact corrected Verlet Gibbs theorem}{Instantiate the exact corrected Verlet Gibbs theorem}}


Take \(\lambda\) to be product Haar phase measure, \(J_h=R\,V_h\) to be the constructed reversed Verlet proposal, and \(w(z)=e^{-H(z)}\). The Hamiltonian module proves that \(J_h\) is both an involution and measure preserving. Measurability of the potential and gradient proves measurability of the weight and proposal, and the exponential weight is strictly positive. The generic theorem therefore proves exact preservation of the actual Gibbs weighted measure by this Metropolis step. Its accepted state is velocity-reversed Verlet and its rejected state is the old state. Applying a further velocity reversal changes those to ordinary Verlet on acceptance and a velocity flip on rejection; any limiting kinetic argument must account for this distinction explicitly.



{\raggedright\noindent Lean declarations: \texttt{metropolis\_\allowbreak{}verlet\_\allowbreak{}preserves\_\allowbreak{}gibbs}.\par}

\end{document}
